\exnumber{05}
\extitle{UX I - Ex 6 - Idéation}
\exfiles{ex05/fichier Figma et PDF, nomme Name\_FirstName\_DeliverableName\_Date\_VersionNumber}
\exauthorize{Ceux nommés dans l'exercice}
\makeheaderfiles

\textbf{Instructions :}

\emph{Vous pouvez remplacer Pierre par votre propre persona.}

L'énoncé de problème choisi est : « How Might We aider Pierre (ou votre persona) à mesurer et préparer son hydratation avant une course ? »

Maintenant, trouvons des solutions créatives ! Souvenez-vous, la solution ne doit pas être numérique ! Oubliez les appareils pour cette quête ;)

Trouvez 3 amis en dehors du programme et organisez une session d'idéation rapide (20 minutes).

\begin{itemize}
\item Ce dont vous avez besoin pour un atelier d'idéation en présentiel :
\begin{itemize}
\item Préparer des papiers ou des post-it, des stickers, un bocal et des stylos.
\item Une salle fermée pendant 20 minutes où tout le monde se sent à l'aise
\end{itemize}
\item Consignes :
\begin{itemize}
\item Demandez-leur d'éteindre leurs téléphones
\item Expliquer le contexte et présenter Pierre (ou votre persona)
\item Expliquer l'objectif :
\begin{itemize}
\item Trouvons autant de solutions que possible à la question : « Comment pourrions-nous aider Pierre (ou votre persona) à mesurer et préparer son hydratation avant une course ? », compte tenu des contraintes suivantes (dévoilez-les une après l'autre : 1 minute par contrainte, 1 contrainte par tour) :
\begin{itemize}
\item Si Pierre était aveugle
\item Si Pierre était un oiseau
\item Si Pierre était milliardaire
\item Si Pierre avait un jumeau constamment avec lui
\item Si Pierre avait 5 ans
\item Si Pierre avait 5 bras
\item Si Pierre avait 105 ans
\item Si Pierre était très maladroit
\item Si Pierre était très oublieux (et perdait la mémoire toutes les 5 minutes)
\item Si Pierre était sourd
\item Si Pierre était une tortue
\item Si Pierre était nu
\item Si Pierre ne pouvait pas s'arrêter de parler 24 h/24 et 7 j/7
\item Si Pierre était toujours en retard
\item Si Pierre était malhonnête
\end{itemize}
\end{itemize}
\item Régler un minuteur pour chaque tour.
\item Demandez aux participants de faire, à chaque tour, ce qui suit :
\begin{itemize}
\item Noter leurs idées sur un post-it (1 idée par post-it)
\item Les lire à voix haute
\item Mettre le post-it dans un bocal
\end{itemize}
\item Rassembler toutes les idées.
\item Donnez 3 stickers à chaque participant et demandez-leur de voter pour les trois meilleures options.
\item S'il n'y a pas de solution qui se démarque clairement, affichez-les sur un tableau et organisez-les selon ces deux paramètres : facile à mettre en œuvre et efficace. Exemple dans les ressources.
\item Choisissez 1 idée principale pour résoudre le problème
\end{itemize}
\end{itemize}

\textbf{Livrables :}

\begin{itemize}
\item Photos de la session
\item Un board avec toutes les idées (sur papier ou en numérique sur Miro ou Figma)
\item Une idée principale pour résoudre le problème à côté du tableau
\item Déposez les documents sur Github, dans un zip nommé « ProjectTitle\_Name\_FirstName », et chaque livrable ainsi : « Name\_FirstName\_DeliverableName\_Date\_VersionNumber ». Exemple : « Doe\_John\_MidFiPrototype\_05242024\_V1 ».
\item Ajoutez un titre dans tous les documents écrits.
\item Soignez l'aspect global des livrables. Ils doivent être clairs, simples et lisibles. Vous pouvez regarder canva.com, sans surcharger le dessin de détails.
\end{itemize}

\textbf{Ressources :}

\begin{itemize}
\item Outil : Framework Ideation Voting Board sur Figma \url{https://www.figma.com/file/nTAWiacCffXvd8dXsLP3Sb/Ideation-Voting-Framework}
\item D'où viennent les bonnes idées par Steven Johnson \url{https://www.ted.com/talks/steven_johnson_where_good_ideas_come_from}
\item Creative Confidence de Tom Keller \url{https://booksvooks.com/creative-confidence-unleashing-the-creative-potential-within-us-all-pdf.html}
\item Podcast What is wrong with UX? \url{https://www.usersknow.com/podcast/2016/10/10/ua1wdsszhx2pemugxc1e0qj0a3l9cs}
\item Brainwriting \#2 \url{https://dux.typepad.com/dux/2011/01/method-2-of-100-brainwriting-brainwriting-is-an-ideation-method-for-quickly-generating-ideas-by-asking-people-to-write-thei.html}
\end{itemize}

\textbf{Citation :}

Il y a de fortes chances que... Si vous avez passé en revue toutes les solutions possibles, vous avez trouvé la bonne.

